\documentclass[../principal.tex]{subfiles}
\graphicspath{{\subfix{../imagenes/}}}

\begin{document}

\chapter{Secuenciadores}

Los secuenciadores son máquinas que generan patrones.

En un sintetizador sonoro se usan para controlar distintos parámetros sonoros, como frecuencia, amplitud, o cualquier otro parámetro que pueda ser controlado por voltaje.

Los secuenciadores presentan los siguientes parámetros posibles:

\begin{enumerate}
    \item Ritmo
    \item Número de pasos
    \item Número de secuencias
    \item Número de canales
\end{enumerate}

\subsection{Secuenciador de 4 pasos}

Existen ejemplos en el ecosistema Eurorack, donde dado el costo y tamaño de cada módulo, han aparecido construcciones de 4 pasos.

En un secuenciador de 4 pasos enumeramos los pasos de 1 a 4.

Con esto después el cambio de paso ocurre al ritmo de un reloj, o con una señal de avanzar.

Cada uno de los pasos tiene asociado un valor del paso, por ejemplo si el secuenciador emite voltajes entre 0 y 5V, puede emitir cualquier valor entre esos extremos.

\subsection{Secuenciador de 8 pasos}

Un ejemplo de estos secuenciadores son los Korg SQ-10, que a su vez son una reversión del antiguo KORG SQ-10.

\newpage

\end{document}

